% translation-source: docs/user/latex/chapters/22-assembler.tex; sha256: b74ee6fa2877dd8f8dcd907849cbcae38b55ff53588a797353fc5bd1438bdafe
\chapter{Assembling several photos}
The \emph{Assemble} tab, at the top of the right panel, makes a new image from the
photos ticked in the ribbon. The original photos do not change. The work runs in the
background and can be cancelled; you can keep browsing your images meanwhile.

\section{Stacking to reduce noise}
Stacking averages several views of the same scene, at the same exposure. Noise falls
with the square root of the number of views: four views halve it, sixteen divide it by
four. Details remain, because they are the same from one view to the next.

\begin{enumerate}
\item Tick at least two photos in the ribbon.
\item Choose the \emph{Reference}: the final image keeps its framing. The middle photo
is proposed, as it stays the closest to the others.
\item Click \emph{Assemble}.
\item Look at the preview, choose the format, then \emph{Export}.
\item \emph{Open the result} brings it into the editor, for further editing.
\end{enumerate}

The views are registered automatically: a tripod is not needed. The result is cropped to
the area common to all views. What moves from one view to the next, like a passer-by, is
left out if it appears in only a minority of views. A view that cannot be registered or
whose exposure differs too much is left out, and the message says so.

\section{Choosing the format}
\begin{itemize}
\item \emph{DNG}: the camera's signal, already corrected for the lens, with no profile
applied. The file reopens like a RAW: profile, balance and settings stay free.
\item \emph{TIFF}: the image rendered in 16 bits with the reference's profile, readable
everywhere.
\end{itemize}
An existing file is never replaced: the name receives a number.

\section{Assembling a panorama}
Choose \emph{Panorama} at the top of the tab. Photos taken from the same place make one wide
view, on one row, horizontal or vertical, or on several rows. Tick them in the ribbon, in any
order, then click \emph{Assemble}.

Let neighbouring photos overlap by a quarter to a third, and keep the same shooting settings.
A photo that overlaps no other, or taken at another focal length, is left out, and the message
says so.

\begin{itemize}
\item \emph{The bottom of the view} is what the panorama leans on: it stays level, and the
crop rather takes from the top.
\item \emph{Projection}: \emph{Automatic} chooses from the field covered. \emph{Flat} keeps
lines straight over a narrow field, \emph{Cylindrical} suits a long sweep, \emph{Spherical} a
field wide in both directions.
\item \emph{Edges}: a panorama seldom has straight edges. \emph{Crop} keeps the largest
rectangle of real pixels; it is the default, and nothing is invented. \emph{Fill} keeps the
whole surface and completes the corners with pixels invented from the rest of the image.
\emph{Leave empty} keeps the empty corners, to crop later in the Geometry tab.
\end{itemize}
Filling is convincing on a sky. On rock or vegetation, look at the result at 100~\%: some blocks
may be out of place. The final message tells the share of invented pixels.

Hand-held, a close foreground moves from one photo to the next: the seams run where it shows
least as far as they can, but turn around the camera rather than around yourself.

\section{Photographing the Moon}
A small subject on an empty background, like the Moon, gets one more treatment. The
deformations of the atmosphere are corrected on the subject, and only the sharpest half
of the views is merged there. The rest of the sky keeps all the views. The message tells
when this correction took place. In our measurements at 200~mm, the choice of the sharp
views brings most of the gain: the longer the series, the more good views to choose from.

\section{Limits}
\begin{itemize}
\item The capture settings must stay the same: stacking is not HDR.
\item A near object shot hand-held may keep a slightly soft edge: it does not move like
the background.
\item The work needs about 400~MB of temporary disk per 32~Mpx photo for a stack, and twice
that for a panorama.
\item A panorama does not warp its edges to fill the frame.
\item A colour profile for the camera is required.
\end{itemize}
